% !PRINTED 555
\weberpage{555}
%<<CONT
In the exponent of $r$ only the residue with respect to the modulus $n$ matters, since $r^h = r^k$ holds always and only when $h \equiv k \pmod{n}$.

As we have already seen in the twelfth section, the quantities (1) are the roots of an \gesperrt{irreducible equation} of degree $(n-1)$
\weq{2}{X = x^{n-1} + x^{n-2} + x^{n-3} + \cdots + x + 1 = 0,}
and we wish first to show that this equation is \gesperrt{cyclic}. This property results from the existence of primitive roots of the prime $n$, with which we occupied ourselves in \S.\,136. It was proved there that for every prime $n$ there are certain numbers $g$, called primitive roots of $n$, which are characterised by the property that among the residues of the powers
\[
1, \ g, \ g^2, \ g^3 \dots g^{n-2},
\]
upon division by $n$, each of the numbers
% [sic] authority includes 0 here; powers of a primitive root modulo n cover only 1,...,n-1.
\[
0, \ 1, \ 2, \dots n-1
\]
occurs once and only once. The residue of a power $g^h$ remains the same when $h$ is changed by a multiple of $n-1$. If
\[
g^\alpha \equiv a \pmod{n},
\]
then $\alpha$ is called the index of $a$ ($\alpha = \operatorname{ind} a$) and $a$ the number or the numerus. The index is taken with respect to the modulus $n-1$, the numerus with respect to the modulus $n$ (\S.\,136).

If, then, we assume such a primitive root of $n$, we can, apart from the order, represent the quantities (1) thus:
\[
r, \ r^g, \ r^{g^2} \dots r^{g^{n-2}},
\]
or, if we set
\weq{3}{r^{g^h} = r_h}
then
\weq{4}{r, \ r_1, \ r_2 \dots r_{n-2}.}
Every number of the field $R(r)$, i.\,e. every rational function of $r$, can be brought into the form
\weq{5}{\varphi(r) = b_0 + b_1 r + b_2 r^2 + \cdots + b_{n-2} r^{n-2},}
or, since by (2)
\[
1 + r + r^2 + \cdots + r^{n-1} = 0
\]
holds,
%CONT>>
